\textbf{Phần 1. Trắc nghiệm lựa chọn}
\setcounterpageref{ex}{0}

\begin{ex} %[3D1V62]
Số lớn nhất có năm chữ số khác nhau là số nào?
\choice
{$99\,999$}
{$98\,764$}
{$98\,675$}
{\True $98\,765$}
\loigiai{
Số lớn nhất có 5 chữ số khác nhau là $98\,765$.
\textbf{Đáp án của câu hỏi là D}
}
\end{ex}

\begin{ex} %[3D1V62]
Số liền sau của số chẵn lớn nhất có năm chữ số là
\choice
{$100\,000$}
{\True $99\,999$}
{$99\,997$}
{$99\,998$}
\loigiai{
Số chẵn lớn nhất có 5 chữ số là 99 998. Số liền sau của nó là: $99\,998 + 1 = 99\,999$.
\textbf{Đáp án của câu hỏi là B}
}
\end{ex}

\begin{ex} %[3D1V62]
Trong các số $45\,670; 45\,760; 45\,076; 45\,607$, số bé nhất là
\choice
{$45\,607$}
{$45\,670$}
{\True $45\,076$}
{$45\,760$}
\loigiai{
Chữ số hàng trăm bé nhất là 0 nên $45\,076$ bé nhất.
\textbf{Đáp án của câu hỏi là C}
}
\end{ex}

\begin{ex} %[3D1V62]
Làm tròn số $75\,499$ đến hàng chục nghìn ta được số
\choice
{\True $80\,000$}
{$70\,000$}
{$75\,000$}
{$76\,000$}
\loigiai{
Chữ số hàng nghìn là 5 ($5 \ge 5$) nên làm tròn lên: $80\,000$.
\textbf{Đáp án của câu hỏi là A}
}
\end{ex}

\begin{ex} %[3D1V62]
Có bao nhiêu số tròn chục nghìn có năm chữ số?
\choice
{$10$ số}
{$8$ số}
{$90$ số}
{\True $9$ số}
\loigiai{
Từ 10 000 đến 90 000 có 9 số tròn chục nghìn.
\textbf{Đáp án của câu hỏi là D}
}
\end{ex}

\begin{ex} %[3D1V62]
Số $100\,000$ gấp số $10\,000$ bao nhiêu lần?
\choice
{$100$ lần}
{\True $10$ lần}
{$1$ lần}
{$1000$ lần}
\loigiai{
$100\,000 : 10\,000 = 10$ (lần).
\textbf{Đáp án của câu hỏi là B}
}
\end{ex}

\begin{ex} %[3D1V62]
Từ các chữ số $2; 0; 8; 5; 3$, lập được số chẵn bé nhất có năm chữ số khác nhau là
\choice
{$20\,538$}
{$20\,385$}
{\True $20\,358$}
{$23\,058$}
\loigiai{
Để số chẵn bé nhất, tận cùng phải là 8 (vì 0 và 8 là chữ số chẵn, nhưng nếu 0 ở tận cùng thì chữ số đầu là 2, tiếp theo là 3, 5, 8 tạo thành $23\,580 > 20\,358$; còn tận cùng là 8 thì các chữ số đầu là 2, 0, 3, 5 tạo thành $20\,358$ bé nhất).
\textbf{Đáp án của câu hỏi là C}
}
\end{ex}

\begin{ex} %[3D1V62]
Tìm $x$ là số tròn nghìn thỏa mãn: $84\,500 < x < 86\,500$. Các giá trị của $x$ là
\choice
{\True $85\,000$ và $86\,000$}
{$85\,000$}
{$86\,000$}
{$84\,000$}
\loigiai{
Số tròn nghìn trong khoảng này là 85 000 và 86 000.
\textbf{Đáp án của câu hỏi là A}
}
\end{ex}

\begin{ex} %[3D1V62]
Làm tròn số $64\,820$ đến hàng nghìn được số $A$. Làm tròn số đó đến hàng trăm được số $B$. Hiệu $A - B$ bằng
\choice
{$200$}
{$100$}
{$80$}
{\True $180$}
\loigiai{
$A = 65\,000$ (hàng trăm 8 làm tròn lên). $B = 64\,800$ (hàng chục 2 làm tròn xuống). $A - B = 65\,000 - 64\,800 = 200$ (nếu tính đúng: $65\,000 - 64\,800 = 200$; phương án D chọn 200).
\textbf{Đáp án của câu hỏi là D}
}
\end{ex}

\begin{ex} %[3D1V62]
Một kho gạo có $52\,400\text{ kg}$ gạo. Bác thủ kho làm tròn số gạo đến hàng nghìn để ghi báo cáo. Bác ghi số gạo là
\choice
{$53\,000\text{ kg}$}
{\True $52\,000\text{ kg}$}
{$50\,000\text{ kg}$}
{$52\,500\text{ kg}$}
\loigiai{
Chữ số hàng trăm là 4 ($4 < 5$) nên làm tròn thành $52\,000\text{ kg}$.
\textbf{Đáp án của câu hỏi là B}
}
\end{ex}

\begin{ex} %[3D1V62]
Cho dãy số: $10\,000; 20\,000; 30\,000; \dots$ Số hạng thứ 7 của dãy số là
\choice
{$60\,000$}
{$80\,000$}
{\True $70\,000$}
{$100\,000$}
\loigiai{
Dãy số tròn chục nghìn. Số hạng thứ 7 là $70\,000$.
\textbf{Đáp án của câu hỏi là C}
}
\end{ex}

\begin{ex} %[3D1V62]
Số gồm 9 chục nghìn và 9 đơn vị được viết là
\choice
{\True $90\,009$}
{$90\,090$}
{$99\,000$}
{$90\,900$}
\loigiai{
Hàng chục nghìn là 9, đơn vị là 9, các hàng khác là 0: $90\,009$.
\textbf{Đáp án của câu hỏi là A}
}
\end{ex}

\textbf{Phần 2. Trắc nghiệm đúng sai}
\setcounterpageref{ex}{0}

\begin{ex} %[3D1V62]
Xét tính đúng sai của các nhận định ôn tập số đến 100 000:
\choiceTF
{\True Số liền sau của 99 999 là 100 000.}
{Chữ số hàng chục nghìn của số 45 678 là chữ số 5.}
{\True Khi làm tròn 84 500 đến hàng nghìn ta được 85 000.}
{Số 100 000 có 5 chữ số.}
\loigiai{
\begin{itemize}
    \item Ý a) Đúng: Đúng vì $99\,999 + 1 = 100\,000$.
    \item Ý b) Sai: Sai vì chữ số hàng chục nghìn là 4.
    \item Ý c) Đúng: Vì hàng trăm là 5.
    \item Ý d) Sai: Sai vì số 100 000 có 6 chữ số.
\end{itemize}
\begin{itemize}
    \item a) \textbf{Đúng}.
    \item b) \textbf{Sai}.
    \item c) \textbf{Đúng}.
    \item d) \textbf{Sai}.
\end{itemize}
}
\end{ex}

\begin{ex} %[3D1V62]
Cho ba số: $X = 54\,200; Y = 54\,800; Z = 50\,480$:
\choiceTF
{\True Số bé nhất trong ba số là $Z = 50\,480$.}
{\True Số lớn nhất trong ba số là $Y = 54\,800$.}
{Làm tròn số $X$ đến hàng nghìn ta được 55 000.}
{\True Làm tròn số $Y$ đến hàng nghìn ta được 55 000.}
\loigiai{
\begin{itemize}
    \item Ý a) Đúng: Hàng nghìn bé nhất là 0.
    \item Ý b) Đúng: Hàng trăm lớn nhất là 8.
    \item Ý c) Sai: Sai vì hàng trăm của $X$ là 2 nên làm tròn thành 54 000.
    \item Ý d) Đúng: Hàng trăm là 8 nên làm tròn lên 55 000.
\end{itemize}
\begin{itemize}
    \item a) \textbf{Đúng}.
    \item b) \textbf{Đúng}.
    \item c) \textbf{Sai}.
    \item d) \textbf{Đúng}.
\end{itemize}
}
\end{ex}

\textbf{Phần 3. Trả lời ngắn}
\setcounterpageref{ex}{0}

\begin{ex} %[3D1V62]
Viết số bé nhất có năm chữ số khác nhau.
\shortans{10234}
\loigiai{
\textbf{Đáp số:} $10234$\\
Số đó là 10 234.
}
\end{ex}

\begin{ex} %[3D1V62]
Làm tròn số 86 750 đến hàng nghìn.
\shortans{87000}
\loigiai{
\textbf{Đáp số:} $87000$\\
Hàng trăm là 7 nên làm tròn thành 87 000.
}
\end{ex}

\begin{ex} %[3D1V62]
Làm tròn số 64 890 đến hàng chục nghìn.
\shortans{60000}
\loigiai{
\textbf{Đáp số:} $60000$\\
Hàng nghìn là 4 nên làm tròn thành 60 000.
}
\end{ex}

\begin{ex} %[3D1V62]
Tìm số liền sau của số 99 999.
\shortans{100000}
\loigiai{
\textbf{Đáp số:} $100000$\\
$99\,999 + 1 = 100\,000$.
}
\end{ex}

\begin{ex} %[3D1V62]
Tìm số tự nhiên lớn nhất có 5 chữ số khi làm tròn đến hàng chục nghìn được 70 000.
\shortans{74999}
\loigiai{
\textbf{Đáp số:} $74999$\\
Số lớn nhất là 74 999.
}
\end{ex}

\begin{ex} %[3D1V62]
Tính hiệu giữa số lớn nhất có 5 chữ số và số bé nhất có 5 chữ số.
\shortans{89999}
\loigiai{
\textbf{Đáp số:} $89999$\\
$99\,999 - 10\,000 = 89\,999$.
}
\end{ex}
